\chapter{Phân tích chức năng AI và cá nhân hóa}
\label{chap:ai-personalization}

Chương này mô tả các chức năng cốt lõi liên quan đến trí tuệ nhân tạo
và cá nhân hóa trong hệ thống. Các chức năng sử dụng dữ liệu từ kết quả
đánh giá, lịch sử luyện tập, hồ sơ năng lực, điểm mục tiêu và tiến độ học
tập để cung cấp nội dung phù hợp cho từng học sinh.

Các kết quả do thành phần AI tạo ra có vai trò hỗ trợ học tập và ra quyết
định. Hệ thống không xem kết quả AI là kết luận tuyệt đối. Trong những
trường hợp cần thiết, giáo viên có thể xem xét và điều chỉnh kết quả.

% =========================================================
\section{UC-08: Chấm điểm và phân tích kết quả làm bài}
\label{sec:uc08-score-analysis}
% =========================================================

\subsection{Thông tin chung}

\renewcommand{\arraystretch}{1.25}

\begin{longtable}{|p{4.2cm}|p{10.2cm}|}
	\caption{Thông tin UC-08: Chấm điểm và phân tích kết quả làm bài}
	\label{tab:uc08-info}\\
	\hline
	\textbf{Thuộc tính} & \textbf{Nội dung} \\
	\hline
	\endfirsthead
	
	\hline
	\textbf{Thuộc tính} & \textbf{Nội dung} \\
	\hline
	\endhead
	
	Mã Use Case &
	UC-08 \\
	\hline
	
	Tên Use Case &
	Chấm điểm và phân tích kết quả làm bài \\
	\hline
	
	Tác nhân chính &
	Student \\
	\hline
	
	Tác nhân phụ &
	Teacher, hệ thống chấm điểm và thành phần AI phân tích kết quả \\
	\hline
	
	Mục đích &
	Tự động chấm điểm bài đánh giá hoặc đề thi thử, phân tích kết quả theo
	từng nhóm kiến thức và cung cấp nhận xét về điểm mạnh, điểm yếu của học
	sinh. \\
	\hline
	
	Tiền điều kiện &
	Student đã đăng nhập; bài đánh giá đang hợp lệ; hệ thống có đáp án,
	thang điểm và thông tin năng lực liên kết với từng câu hỏi. \\
	\hline
	
	Hậu điều kiện &
	Kết quả được lưu; báo cáo phân tích được tạo; dữ liệu được chuyển sang
	chức năng cập nhật hồ sơ năng lực. \\
	\hline
	
	Tần suất sử dụng &
	Sau mỗi lần Student nộp bài đánh giá, bài luyện tập hoặc đề thi thử. \\
	\hline
	
\end{longtable}

\subsection{Luồng xử lý chính}

\begin{enumerate}
	\item Student hoàn thành bài đánh giá hoặc đề thi thử.
	
	\item Student chọn chức năng nộp bài.
	
	\item Hệ thống kiểm tra trạng thái bài làm và tính hợp lệ của dữ liệu.
	
	\item Hệ thống ghi nhận thời điểm nộp bài và thời gian làm bài.
	
	\item Hệ thống đối chiếu câu trả lời của Student với đáp án.
	
	\item Hệ thống tính điểm cho từng câu hỏi.
	
	\item Hệ thống tính điểm tổng và điểm theo từng phần của bài thi.
	
	\item Hệ thống thống kê số câu đúng, số câu sai, số câu chưa trả lời
	và thời gian trung bình cho mỗi câu.
	
	\item Hệ thống nhóm kết quả theo môn học, chủ đề và năng lực.
	
	\item Thành phần phân tích xác định các nhóm kiến thức Student đang
	thực hiện tốt.
	
	\item Thành phần phân tích xác định các nhóm kiến thức Student còn yếu.
	
	\item Hệ thống tạo nhận xét và đề xuất nội dung cần ưu tiên ôn tập.
	
	\item Hệ thống lưu kết quả bài làm và báo cáo phân tích.
	
	\item Hệ thống gửi dữ liệu sang chức năng cập nhật hồ sơ năng lực.
	
	\item Hệ thống hiển thị kết quả cho Student.
\end{enumerate}

\subsection{Luồng thay thế và trường hợp lỗi}

\begin{enumerate}[label=\textbf{A\arabic*.}]
	\item \textbf{Bài làm còn câu chưa trả lời}
	
	Hệ thống hiển thị cảnh báo cho Student. Student có thể quay lại làm
	tiếp hoặc xác nhận nộp bài với các câu chưa trả lời.
	
	\item \textbf{Bài có câu hỏi tự luận}
	
	Hệ thống tự động chấm các câu hỏi trắc nghiệm. Các câu tự luận được
	chuyển cho Teacher chấm trước khi hoàn thiện điểm tổng.
	
	\item \textbf{Câu hỏi không có đáp án hợp lệ}
	
	Hệ thống tạm thời không tính điểm câu hỏi đó, ghi nhận lỗi dữ liệu và
	thông báo cho người quản lý ngân hàng câu hỏi.
	
	\item \textbf{Mất kết nối khi nộp bài}
	
	Hệ thống giữ bản lưu gần nhất. Student có thể gửi lại khi kết nối
	được khôi phục.
	
	\item \textbf{Thành phần AI phân tích không hoạt động}
	
	Hệ thống vẫn thực hiện chấm điểm theo đáp án và hiển thị thống kê cơ
	bản. Phần nhận xét AI được đánh dấu là tạm thời chưa khả dụng.
	
	\item \textbf{Bài làm đã hết thời gian}
	
	Hệ thống tự động nộp bài tại thời điểm hết giờ và tiếp tục thực hiện
	quá trình chấm điểm.
\end{enumerate}

\subsection{Dữ liệu đầu vào và kết quả đầu ra}

\begin{longtable}{|p{3.5cm}|p{5.3cm}|p{5.6cm}|}
	\caption{Dữ liệu của UC-08}
	\label{tab:uc08-data}\\
	\hline
	\textbf{Nhóm dữ liệu} &
	\textbf{Thành phần} &
	\textbf{Mô tả} \\
	\hline
	\endfirsthead
	
	\hline
	\textbf{Nhóm dữ liệu} &
	\textbf{Thành phần} &
	\textbf{Mô tả} \\
	\hline
	\endhead
	
	Đầu vào &
	Câu trả lời của Student &
	Danh sách đáp án Student đã chọn hoặc nội dung Student đã nhập. \\
	\hline
	
	Đầu vào &
	Đáp án và thang điểm &
	Đáp án đúng, trọng số và quy tắc tính điểm của bài đánh giá. \\
	\hline
	
	Đầu vào &
	Thông tin câu hỏi &
	Môn học, chủ đề, năng lực, độ khó và loại câu hỏi. \\
	\hline
	
	Đầu vào &
	Thời gian làm bài &
	Thời điểm bắt đầu, kết thúc và thời gian cho từng câu hỏi. \\
	\hline
	
	Đầu ra &
	Điểm tổng và điểm thành phần &
	Điểm toàn bài và điểm theo từng môn học hoặc nhóm năng lực. \\
	\hline
	
	Đầu ra &
	Thống kê bài làm &
	Số câu đúng, sai, bỏ trống và thời gian làm bài. \\
	\hline
	
	Đầu ra &
	Phân tích năng lực &
	Danh sách điểm mạnh, điểm yếu và chủ đề cần ưu tiên. \\
	\hline
	
	Đầu ra &
	Gợi ý ôn tập &
	Nội dung, tài liệu hoặc bài luyện tập được đề xuất. \\
	\hline
	
\end{longtable}

\subsection{Mối liên hệ với hồ sơ năng lực và lộ trình học}

Kết quả của UC-08 là nguồn dữ liệu chính cho UC-09. Điểm theo từng nhóm
kiến thức được sử dụng để cập nhật mức độ thành thạo trong hồ sơ năng
lực.

Khi hồ sơ năng lực thay đổi, UC-11 có thể tạo mới hoặc điều chỉnh lộ
trình học cá nhân hóa. Những nhóm kiến thức có kết quả thấp được tăng
mức độ ưu tiên trong lộ trình.

% =========================================================
\section{UC-09: Tạo hoặc cập nhật hồ sơ năng lực học sinh}
\label{sec:uc09-learning-profile}
% =========================================================

\subsection{Thông tin chung}

\begin{longtable}{|p{4.2cm}|p{10.2cm}|}
	\caption{Thông tin UC-09: Tạo hoặc cập nhật hồ sơ năng lực}
	\label{tab:uc09-info}\\
	\hline
	\textbf{Thuộc tính} & \textbf{Nội dung} \\
	\hline
	\endfirsthead
	
	\hline
	\textbf{Thuộc tính} & \textbf{Nội dung} \\
	\hline
	\endhead
	
	Mã Use Case &
	UC-09 \\
	\hline
	
	Tên Use Case &
	Tạo hoặc cập nhật hồ sơ năng lực học sinh \\
	\hline
	
	Tác nhân chính &
	Hệ thống \\
	\hline
	
	Tác nhân liên quan &
	Student và Teacher \\
	\hline
	
	Mục đích &
	Tổng hợp dữ liệu học tập để biểu diễn mức độ thành thạo, điểm mạnh,
	điểm yếu và xu hướng tiến bộ của Student. \\
	\hline
	
	Tiền điều kiện &
	Student có tài khoản hợp lệ và có ít nhất một kết quả đánh giá, luyện
	tập hoặc thi thử đã được lưu. \\
	\hline
	
	Hậu điều kiện &
	Hồ sơ năng lực mới được tạo hoặc phiên bản hiện tại được cập nhật và lưu
	trong hệ thống. \\
	\hline
	
	Tần suất sử dụng &
	Sau khi có kết quả học tập mới hoặc theo chu kỳ cập nhật của hệ thống. \\
	\hline
	
\end{longtable}

\subsection{Luồng xử lý chính}

\begin{enumerate}
	\item Hệ thống nhận kết quả mới từ UC-08 hoặc UC-12.
	
	\item Hệ thống xác định Student liên quan đến kết quả.
	
	\item Hệ thống truy xuất lịch sử đánh giá, luyện tập và thi thử.
	
	\item Hệ thống nhóm dữ liệu theo môn học, chủ đề và năng lực.
	
	\item Hệ thống tính mức độ thành thạo hiện tại của từng năng lực.
	
	\item Hệ thống so sánh kết quả mới với các kết quả trước đó.
	
	\item Hệ thống xác định xu hướng tăng, giảm hoặc ổn định.
	
	\item Hệ thống xác định các nhóm kiến thức mạnh.
	
	\item Hệ thống xác định các khoảng trống kiến thức hoặc kỹ năng yếu.
	
	\item Hệ thống tạo mới hoặc cập nhật hồ sơ năng lực.
	
	\item Hệ thống lưu thời điểm cập nhật và phiên bản hồ sơ.
	
	\item Hệ thống thông báo cho UC-11 rằng hồ sơ năng lực đã thay đổi.
	
	\item Student và Teacher có thể xem hồ sơ năng lực mới nhất.
\end{enumerate}

\subsection{Luồng thay thế và trường hợp lỗi}

\begin{enumerate}[label=\textbf{A\arabic*.}]
	\item \textbf{Student chưa có đủ dữ liệu}
	
	Hệ thống tạo hồ sơ tạm thời và hiển thị cảnh báo rằng mức độ chính
	xác còn hạn chế.
	
	\item \textbf{Kết quả mới có giá trị bất thường}
	
	Hệ thống giữ lại dữ liệu nhưng đánh dấu để Teacher kiểm tra trước khi
	sử dụng cho quyết định cá nhân hóa.
	
	\item \textbf{Kết quả bài làm bị hủy}
	
	Hệ thống loại bỏ tác động của kết quả đó và tính toán lại hồ sơ.
	
	\item \textbf{Thiếu thông tin liên kết năng lực của câu hỏi}
	
	Hệ thống vẫn lưu điểm bài làm nhưng chưa sử dụng phần dữ liệu đó để
	cập nhật năng lực cụ thể.
	
	\item \textbf{Quá trình tính toán bị lỗi}
	
	Hệ thống giữ nguyên phiên bản hồ sơ gần nhất và ghi nhận lỗi để xử lý.
\end{enumerate}

\subsection{Dữ liệu đầu vào và kết quả đầu ra}

\begin{longtable}{|p{3.5cm}|p{5.3cm}|p{5.6cm}|}
	\caption{Dữ liệu của UC-09}
	\label{tab:uc09-data}\\
	\hline
	\textbf{Nhóm dữ liệu} &
	\textbf{Thành phần} &
	\textbf{Mô tả} \\
	\hline
	\endfirsthead
	
	\hline
	\textbf{Nhóm dữ liệu} &
	\textbf{Thành phần} &
	\textbf{Mô tả} \\
	\hline
	\endhead
	
	Đầu vào &
	Kết quả đánh giá &
	Điểm tổng, điểm thành phần và kết quả theo từng năng lực. \\
	\hline
	
	Đầu vào &
	Lịch sử luyện tập &
	Tỷ lệ đúng, độ khó, chủ đề và số lần luyện tập. \\
	\hline
	
	Đầu vào &
	Kết quả thi thử &
	Điểm thi thử, thời gian làm và xu hướng qua các lần thi. \\
	\hline
	
	Đầu vào &
	Tiến độ học tập &
	Tỷ lệ hoàn thành bài học, lộ trình và nhiệm vụ học tập. \\
	\hline
	
	Đầu ra &
	Mức độ thành thạo &
	Giá trị thể hiện năng lực hiện tại của Student theo từng nhóm. \\
	\hline
	
	Đầu ra &
	Điểm mạnh và điểm yếu &
	Danh sách năng lực tốt và năng lực cần cải thiện. \\
	\hline
	
	Đầu ra &
	Xu hướng học tập &
	Mức tăng, giảm hoặc ổn định theo từng giai đoạn. \\
	\hline
	
	Đầu ra &
	Khoảng trống kiến thức &
	Các chủ đề chưa đạt yêu cầu hoặc chưa có đủ dữ liệu. \\
	\hline
	
\end{longtable}

\subsection{Mối liên hệ với lộ trình học}

Hồ sơ năng lực là dữ liệu đầu vào quan trọng nhất của UC-11. Khi mức độ
thành thạo thay đổi, hệ thống phải đánh giá lại thứ tự ưu tiên của nội
dung trong lộ trình.

UC-12 sử dụng hồ sơ năng lực để lựa chọn chủ đề và độ khó của câu hỏi.
Sau mỗi phiên luyện tập, kết quả mới tiếp tục quay lại cập nhật hồ sơ,
tạo thành vòng lặp cá nhân hóa liên tục.

% =========================================================
\section{UC-11: Tạo hoặc cập nhật lộ trình học cá nhân hóa}
\label{sec:uc11-learning-path}
% =========================================================

\subsection{Thông tin chung}

\begin{longtable}{|p{4.2cm}|p{10.2cm}|}
	\caption{Thông tin UC-11: Tạo hoặc cập nhật lộ trình học cá nhân hóa}
	\label{tab:uc11-info}\\
	\hline
	\textbf{Thuộc tính} & \textbf{Nội dung} \\
	\hline
	\endfirsthead
	
	\hline
	\textbf{Thuộc tính} & \textbf{Nội dung} \\
	\hline
	\endhead
	
	Mã Use Case &
	UC-11 \\
	\hline
	
	Tên Use Case &
	Tạo hoặc cập nhật lộ trình học cá nhân hóa \\
	\hline
	
	Tác nhân chính &
	Student \\
	\hline
	
	Tác nhân phụ &
	Hệ thống cá nhân hóa và Teacher \\
	\hline
	
	Mục đích &
	Tạo kế hoạch học tập phù hợp với năng lực hiện tại, điểm mục tiêu, thời
	gian còn lại và khả năng học tập của Student. \\
	\hline
	
	Tiền điều kiện &
	Student đã thiết lập điểm mục tiêu; có hồ sơ năng lực; hệ thống có tài
	liệu và bài luyện tập phù hợp. \\
	\hline
	
	Hậu điều kiện &
	Một lộ trình mới được tạo hoặc lộ trình hiện tại được cập nhật và lưu. \\
	\hline
	
	Tần suất sử dụng &
	Khi Student yêu cầu, khi hồ sơ năng lực thay đổi đáng kể hoặc khi tiến
	độ thực tế chênh lệch so với kế hoạch. \\
	\hline
	
\end{longtable}

\subsection{Luồng xử lý chính}

\begin{enumerate}
	\item Student mở chức năng lộ trình học cá nhân hóa.
	
	\item Student yêu cầu tạo mới hoặc cập nhật lộ trình.
	
	\item Hệ thống lấy hồ sơ năng lực mới nhất.
	
	\item Hệ thống lấy điểm mục tiêu của Student.
	
	\item Hệ thống lấy ngày thi dự kiến và thời gian còn lại.
	
	\item Hệ thống lấy thời gian học trung bình mỗi ngày hoặc mỗi tuần.
	
	\item Hệ thống xác định khoảng cách giữa năng lực hiện tại và mục tiêu.
	
	\item Hệ thống xếp hạng các môn học, chủ đề và năng lực cần ưu tiên.
	
	\item Hệ thống lựa chọn tài liệu, bài học và bài luyện tập phù hợp.
	
	\item Hệ thống chia kế hoạch theo tuần hoặc theo giai đoạn.
	
	\item Hệ thống xác định các mốc đánh giá lại năng lực.
	
	\item Hệ thống tạo lộ trình học dự kiến.
	
	\item Hệ thống hiển thị lộ trình cho Student.
	
	\item Student xem và xác nhận thời gian học.
	
	\item Hệ thống lưu lộ trình.
	
	\item Teacher có thể xem và đưa ra đề xuất điều chỉnh.
\end{enumerate}

\subsection{Luồng thay thế và trường hợp lỗi}

\begin{enumerate}[label=\textbf{A\arabic*.}]
	\item \textbf{Student chưa thiết lập điểm mục tiêu}
	
	Hệ thống yêu cầu Student thiết lập điểm mục tiêu trước khi tạo lộ
	trình.
	
	\item \textbf{Chưa có hồ sơ năng lực}
	
	Hệ thống yêu cầu Student thực hiện bài đánh giá đầu vào.
	
	\item \textbf{Không đủ thời gian để hoàn thành toàn bộ nội dung}
	
	Hệ thống tạo lộ trình rút gọn, ưu tiên các chủ đề có mức ảnh hưởng cao
	đến điểm mục tiêu.
	
	\item \textbf{Không có tài liệu phù hợp}
	
	Hệ thống ghi nhận chủ đề còn thiếu tài liệu và đề xuất Teacher bổ
	sung nội dung.
	
	\item \textbf{Student không theo kịp kế hoạch}
	
	Hệ thống điều chỉnh số lượng nhiệm vụ, thời hạn và mức độ ưu tiên.
	
	\item \textbf{Student tiến bộ nhanh hơn kế hoạch}
	
	Hệ thống rút ngắn nội dung đã thành thạo và chuyển sang chủ đề có độ
	khó cao hơn.
	
	\item \textbf{Teacher yêu cầu điều chỉnh}
	
	Hệ thống lưu đề xuất của Teacher và cập nhật lộ trình sau khi được
	xác nhận.
\end{enumerate}

\subsection{Dữ liệu đầu vào và kết quả đầu ra}

\begin{longtable}{|p{3.5cm}|p{5.3cm}|p{5.6cm}|}
	\caption{Dữ liệu của UC-11}
	\label{tab:uc11-data}\\
	\hline
	\textbf{Nhóm dữ liệu} &
	\textbf{Thành phần} &
	\textbf{Mô tả} \\
	\hline
	\endfirsthead
	
	\hline
	\textbf{Nhóm dữ liệu} &
	\textbf{Thành phần} &
	\textbf{Mô tả} \\
	\hline
	\endhead
	
	Đầu vào &
	Hồ sơ năng lực &
	Mức độ thành thạo, điểm mạnh, điểm yếu và khoảng trống kiến thức. \\
	\hline
	
	Đầu vào &
	Điểm mục tiêu &
	Mức điểm Student mong muốn đạt được. \\
	\hline
	
	Đầu vào &
	Ngày thi dự kiến &
	Mốc thời gian kết thúc của kế hoạch học tập. \\
	\hline
	
	Đầu vào &
	Khả năng học tập &
	Thời gian Student có thể dành cho việc học mỗi ngày hoặc mỗi tuần. \\
	\hline
	
	Đầu vào &
	Kho tài liệu &
	Danh sách bài học, tài liệu và bài luyện tập hiện có. \\
	\hline
	
	Đầu ra &
	Kế hoạch theo giai đoạn &
	Danh sách nội dung cần hoàn thành theo tuần hoặc theo mốc. \\
	\hline
	
	Đầu ra &
	Thứ tự ưu tiên &
	Mức độ ưu tiên của từng môn học, chủ đề và năng lực. \\
	\hline
	
	Đầu ra &
	Nhiệm vụ học tập &
	Bài học, tài liệu và bài luyện tập được đề xuất. \\
	\hline
	
	Đầu ra &
	Mốc đánh giá lại &
	Thời điểm Student cần thực hiện bài đánh giá hoặc thi thử. \\
	\hline
	
\end{longtable}

\subsection{Mối liên hệ với hồ sơ năng lực}

Lộ trình được tạo dựa trên hồ sơ năng lực tại thời điểm gần nhất. Khi
UC-09 cập nhật hồ sơ, hệ thống so sánh dữ liệu mới với dữ liệu được dùng
để tạo lộ trình.

Nếu sự thay đổi vượt quá ngưỡng quy định, hệ thống đề xuất cập nhật lộ
trình. Việc này giúp lộ trình không bị cố định và có thể phản ánh đúng
năng lực thực tế của Student.

% =========================================================
\section{UC-12: Luyện tập thích ứng theo năng lực}
\label{sec:uc12-adaptive-practice}
% =========================================================

\subsection{Thông tin chung}

\begin{longtable}{|p{4.2cm}|p{10.2cm}|}
	\caption{Thông tin UC-12: Luyện tập thích ứng theo năng lực}
	\label{tab:uc12-info}\\
	\hline
	\textbf{Thuộc tính} & \textbf{Nội dung} \\
	\hline
	\endfirsthead
	
	\hline
	\textbf{Thuộc tính} & \textbf{Nội dung} \\
	\hline
	\endhead
	
	Mã Use Case &
	UC-12 \\
	\hline
	
	Tên Use Case &
	Luyện tập thích ứng theo năng lực \\
	\hline
	
	Tác nhân chính &
	Student \\
	\hline
	
	Tác nhân phụ &
	Hệ thống lựa chọn câu hỏi và thành phần cá nhân hóa \\
	\hline
	
	Mục đích &
	Lựa chọn câu hỏi và điều chỉnh độ khó theo kết quả thực tế của Student
	trong từng phiên luyện tập. \\
	\hline
	
	Tiền điều kiện &
	Student đã đăng nhập; ngân hàng câu hỏi có dữ liệu; Student có hồ sơ năng
	lực hoặc lịch sử luyện tập. \\
	\hline
	
	Hậu điều kiện &
	Kết quả phiên luyện tập được lưu và hồ sơ năng lực được cập nhật. \\
	\hline
	
	Tần suất sử dụng &
	Theo nhu cầu luyện tập của Student hoặc theo nhiệm vụ trong lộ trình. \\
	\hline
	
\end{longtable}

\subsection{Luồng xử lý chính}

\begin{enumerate}
	\item Student mở chức năng luyện tập thích ứng.
	
	\item Student chọn môn học hoặc chủ đề muốn luyện tập.
	
	\item Hệ thống lấy hồ sơ năng lực và lịch sử luyện tập.
	
	\item Hệ thống xác định mức độ khó khởi đầu phù hợp.
	
	\item Hệ thống chọn một câu hỏi từ ngân hàng câu hỏi.
	
	\item Hệ thống hiển thị câu hỏi cho Student.
	
	\item Student nhập hoặc lựa chọn câu trả lời.
	
	\item Hệ thống chấm câu trả lời.
	
	\item Hệ thống hiển thị phản hồi đúng hoặc sai.
	
	\item Hệ thống cập nhật mức độ thành thạo tạm thời của chủ đề.
	
	\item Hệ thống điều chỉnh độ khó cho câu hỏi tiếp theo.
	
	\item Quy trình tiếp tục cho đến khi Student kết thúc phiên hoặc đạt
	số lượng câu hỏi quy định.
	
	\item Hệ thống tính điểm phiên luyện tập.
	
	\item Hệ thống tạo báo cáo kết quả.
	
	\item Hệ thống gửi dữ liệu sang UC-09 để cập nhật hồ sơ năng lực.
	
	\item Hệ thống cập nhật tiến độ liên quan trong UC-11.
\end{enumerate}

\subsection{Quy tắc điều chỉnh}

\begin{itemize}
	\item Khi Student trả lời đúng nhiều câu liên tiếp, hệ thống tăng độ
	khó hoặc chuyển sang nội dung nâng cao hơn.
	
	\item Khi Student trả lời sai nhiều câu liên tiếp, hệ thống giảm độ
	khó và ưu tiên các câu hỏi củng cố kiến thức nền tảng.
	
	\item Khi Student mất quá nhiều thời gian, hệ thống có thể cung cấp
	gợi ý hoặc chọn câu hỏi có mức độ tương đương nhưng cách diễn đạt
	khác.
	
	\item Chủ đề có mức độ thành thạo thấp được lựa chọn với tần suất cao
	hơn.
	
	\item Hệ thống không lặp lại một câu hỏi quá nhiều lần trong cùng một
	phiên.
\end{itemize}

\subsection{Luồng thay thế và trường hợp lỗi}

\begin{enumerate}[label=\textbf{A\arabic*.}]
	\item \textbf{Không có hồ sơ năng lực}
	
	Hệ thống sử dụng mức độ khó mặc định hoặc yêu cầu Student thực hiện
	bài đánh giá đầu vào.
	
	\item \textbf{Không còn câu hỏi phù hợp}
	
	Hệ thống chuyển sang chủ đề gần nhất hoặc thông báo đã hoàn thành nội
	dung hiện có.
	
	\item \textbf{Student yêu cầu gợi ý}
	
	Hệ thống hiển thị gợi ý theo từng mức nhưng không hiển thị ngay đáp
	án cuối cùng.
	
	\item \textbf{Student thoát giữa phiên}
	
	Hệ thống lưu tiến độ hiện tại và cho phép Student tiếp tục sau.
	
	\item \textbf{Mất kết nối}
	
	Hệ thống lưu câu trả lời gần nhất trên thiết bị hoặc khôi phục từ bản
	lưu gần nhất khi kết nối trở lại.
	
	\item \textbf{Câu hỏi có lỗi}
	
	Student có thể báo lỗi. Hệ thống bỏ qua câu hỏi đó và ghi nhận để
	Teacher kiểm tra.
\end{enumerate}

\subsection{Dữ liệu đầu vào và kết quả đầu ra}

\begin{longtable}{|p{3.5cm}|p{5.3cm}|p{5.6cm}|}
	\caption{Dữ liệu của UC-12}
	\label{tab:uc12-data}\\
	\hline
	\textbf{Nhóm dữ liệu} &
	\textbf{Thành phần} &
	\textbf{Mô tả} \\
	\hline
	\endfirsthead
	
	\hline
	\textbf{Nhóm dữ liệu} &
	\textbf{Thành phần} &
	\textbf{Mô tả} \\
	\hline
	\endhead
	
	Đầu vào &
	Hồ sơ năng lực &
	Mức độ thành thạo của Student theo từng chủ đề. \\
	\hline
	
	Đầu vào &
	Ngân hàng câu hỏi &
	Nội dung, đáp án, chủ đề, độ khó và năng lực liên quan. \\
	\hline
	
	Đầu vào &
	Lịch sử trả lời &
	Các câu đã làm, kết quả và thời gian trả lời. \\
	\hline
	
	Đầu ra &
	Câu hỏi được lựa chọn &
	Câu hỏi phù hợp với năng lực hiện tại. \\
	\hline
	
	Đầu ra &
	Phản hồi tức thời &
	Kết quả đúng hoặc sai, lời giải và gợi ý. \\
	\hline
	
	Đầu ra &
	Điểm phiên luyện tập &
	Điểm số, tỷ lệ đúng và thời gian hoàn thành. \\
	\hline
	
	Đầu ra &
	Mức độ thành thạo mới &
	Dữ liệu được sử dụng để cập nhật hồ sơ năng lực. \\
	\hline
	
\end{longtable}

\subsection{Mối liên hệ với hồ sơ năng lực và lộ trình}

UC-12 sử dụng hồ sơ năng lực để lựa chọn câu hỏi ban đầu. Kết quả của
từng phiên luyện tập được chuyển ngược về UC-09 để cập nhật mức độ thành
thạo.

UC-11 sử dụng tiến độ luyện tập để xác định Student có đang theo kịp lộ
trình hay không. Khi Student chưa đạt yêu cầu, lộ trình có thể bổ sung
thêm bài luyện tập cho chủ đề tương ứng.

% =========================================================
\section{UC-14: Tương tác với AI Tutor}
\label{sec:uc14-ai-tutor}
% =========================================================

\subsection{Thông tin chung}

\begin{longtable}{|p{4.2cm}|p{10.2cm}|}
	\caption{Thông tin UC-14: Tương tác với AI Tutor}
	\label{tab:uc14-info}\\
	\hline
	\textbf{Thuộc tính} & \textbf{Nội dung} \\
	\hline
	\endfirsthead
	
	\hline
	\textbf{Thuộc tính} & \textbf{Nội dung} \\
	\hline
	\endhead
	
	Mã Use Case &
	UC-14 \\
	\hline
	
	Tên Use Case &
	Tương tác với AI Tutor \\
	\hline
	
	Tác nhân chính &
	Student \\
	\hline
	
	Tác nhân phụ &
	AI Tutor và kho tài liệu học tập \\
	\hline
	
	Mục đích &
	Hỗ trợ Student giải thích kiến thức, hướng dẫn cách giải, cung cấp ví dụ
	và gợi ý học tập dựa trên tài liệu của hệ thống. \\
	\hline
	
	Tiền điều kiện &
	Student đã đăng nhập; dịch vụ AI Tutor đang hoạt động; kho tài liệu đã
	được cấu hình. \\
	\hline
	
	Hậu điều kiện &
	Câu trả lời được hiển thị; lịch sử hội thoại và phản hồi của Student
	được lưu. \\
	\hline
	
	Tần suất sử dụng &
	Theo nhu cầu hỏi đáp trong quá trình học tập. \\
	\hline
	
\end{longtable}

\subsection{Luồng xử lý chính}

\begin{enumerate}
	\item Student mở màn hình AI Tutor.
	
	\item Student nhập câu hỏi hoặc yêu cầu giải thích.
	
	\item Hệ thống kiểm tra nội dung đầu vào.
	
	\item Hệ thống xác định môn học, chủ đề và ý định của câu hỏi.
	
	\item Hệ thống lấy ngữ cảnh từ lịch sử hội thoại hiện tại.
	
	\item Hệ thống truy xuất các tài liệu liên quan trong kho dữ liệu.
	
	\item Hệ thống lựa chọn các đoạn nội dung phù hợp.
	
	\item AI Tutor tạo câu trả lời dựa trên câu hỏi và tài liệu được truy
	xuất.
	
	\item Hệ thống kiểm tra định dạng và mức độ phù hợp của câu trả lời.
	
	\item Hệ thống hiển thị câu trả lời cho Student.
	
	\item Hệ thống hiển thị nguồn tài liệu tham khảo nếu có.
	
	\item Student có thể yêu cầu giải thích thêm, đưa ví dụ hoặc trình bày
	theo cách đơn giản hơn.
	
	\item Hệ thống lưu lịch sử hội thoại.
	
	\item Student có thể đánh giá câu trả lời.
\end{enumerate}

\subsection{Nguyên tắc phản hồi của AI Tutor}

\begin{itemize}
	\item Ưu tiên nội dung trong kho tài liệu của hệ thống.
	
	\item Không khẳng định chắc chắn khi không có đủ dữ liệu.
	
	\item Khuyến khích Student hiểu cách giải thay vì chỉ đưa đáp án.
	
	\item Với bài tập, AI Tutor có thể đưa gợi ý theo từng bước.
	
	\item Câu trả lời phải phù hợp với trình độ và hồ sơ năng lực của
	Student.
	
	\item Không cung cấp nội dung ngoài phạm vi học tập được hệ thống cho
	phép.
	
	\item Khi nội dung có độ tin cậy thấp, hệ thống phải hiển thị cảnh báo.
	
	\item AI Tutor không thay thế vai trò chuyên môn của Teacher.
\end{itemize}

\subsection{Luồng thay thế và trường hợp lỗi}

\begin{enumerate}[label=\textbf{A\arabic*.}]
	\item \textbf{Câu hỏi không rõ ràng}
	
	AI Tutor yêu cầu Student cung cấp thêm thông tin hoặc chọn chủ đề.
	
	\item \textbf{Không tìm được tài liệu phù hợp}
	
	Hệ thống thông báo chưa có đủ dữ liệu và đề xuất các chủ đề liên quan.
	
	\item \textbf{Câu hỏi ngoài phạm vi}
	
	Hệ thống từ chối trả lời và hướng Student quay lại nội dung học tập.
	
	\item \textbf{Câu hỏi chứa nội dung không phù hợp}
	
	Hệ thống chặn yêu cầu và ghi nhận sự kiện theo quy định an toàn.
	
	\item \textbf{Dịch vụ AI không hoạt động}
	
	Hệ thống thông báo Student thử lại sau và không làm mất nội dung câu
	hỏi đã nhập.
	
	\item \textbf{Câu trả lời có độ tin cậy thấp}
	
	Hệ thống hiển thị cảnh báo và đề nghị Student kiểm tra lại với tài
	liệu hoặc Teacher.
	
	\item \textbf{Student đánh giá câu trả lời không hữu ích}
	
	Hệ thống lưu phản hồi để phục vụ cải thiện chất lượng AI Tutor.
\end{enumerate}

\subsection{Dữ liệu đầu vào và kết quả đầu ra}

\begin{longtable}{|p{3.5cm}|p{5.3cm}|p{5.6cm}|}
	\caption{Dữ liệu của UC-14}
	\label{tab:uc14-data}\\
	\hline
	\textbf{Nhóm dữ liệu} &
	\textbf{Thành phần} &
	\textbf{Mô tả} \\
	\hline
	\endfirsthead
	
	\hline
	\textbf{Nhóm dữ liệu} &
	\textbf{Thành phần} &
	\textbf{Mô tả} \\
	\hline
	\endhead
	
	Đầu vào &
	Câu hỏi của Student &
	Nội dung Student muốn được giải thích hoặc hướng dẫn. \\
	\hline
	
	Đầu vào &
	Lịch sử hội thoại &
	Các câu hỏi và câu trả lời trước đó trong phiên. \\
	\hline
	
	Đầu vào &
	Hồ sơ năng lực &
	Thông tin giúp điều chỉnh mức độ giải thích cho Student. \\
	\hline
	
	Đầu vào &
	Kho tài liệu &
	Tài liệu học tập được sử dụng làm nguồn tham khảo. \\
	\hline
	
	Đầu ra &
	Câu trả lời &
	Nội dung giải thích, hướng dẫn hoặc gợi ý. \\
	\hline
	
	Đầu ra &
	Ví dụ và bước giải &
	Ví dụ minh họa hoặc các bước hỗ trợ Student tự giải. \\
	\hline
	
	Đầu ra &
	Nguồn tham khảo &
	Tài liệu hoặc nội dung liên quan được sử dụng. \\
	\hline
	
	Đầu ra &
	Lịch sử hội thoại &
	Dữ liệu phục vụ tiếp tục tương tác và đánh giá chất lượng. \\
	\hline
	
\end{longtable}

\subsection{Mối liên hệ với hồ sơ năng lực và lộ trình}

AI Tutor có thể sử dụng hồ sơ năng lực để điều chỉnh mức độ giải thích.
Với Student đang yếu một chủ đề, AI Tutor ưu tiên cách trình bày cơ bản,
ví dụ đơn giản và gợi ý từng bước.

Các chủ đề Student hỏi nhiều hoặc đánh giá là khó hiểu có thể được ghi
nhận làm tín hiệu bổ sung cho hồ sơ năng lực. UC-11 có thể đề xuất thêm
tài liệu hoặc bài luyện tập liên quan đến các chủ đề đó.

% =========================================================
\section{UC-17: Dự đoán khả năng đạt điểm mục tiêu}
\label{sec:uc17-score-prediction}
% =========================================================

\subsection{Thông tin chung}

\begin{longtable}{|p{4.2cm}|p{10.2cm}|}
	\caption{Thông tin UC-17: Dự đoán khả năng đạt điểm mục tiêu}
	\label{tab:uc17-info}\\
	\hline
	\textbf{Thuộc tính} & \textbf{Nội dung} \\
	\hline
	\endfirsthead
	
	\hline
	\textbf{Thuộc tính} & \textbf{Nội dung} \\
	\hline
	\endhead
	
	Mã Use Case &
	UC-17 \\
	\hline
	
	Tên Use Case &
	Dự đoán khả năng đạt điểm mục tiêu \\
	\hline
	
	Tác nhân chính &
	Student \\
	\hline
	
	Tác nhân liên quan &
	Teacher và thành phần dự đoán \\
	\hline
	
	Mục đích &
	Ước lượng khả năng Student đạt điểm mục tiêu dựa trên năng lực hiện tại,
	kết quả thi thử, tiến độ và hành vi học tập. \\
	\hline
	
	Tiền điều kiện &
	Student đã thiết lập điểm mục tiêu; có dữ liệu học tập đủ để phân tích;
	thành phần dự đoán đang hoạt động. \\
	\hline
	
	Hậu điều kiện &
	Kết quả dự đoán, mức độ tin cậy và gợi ý cải thiện được hiển thị và lưu. \\
	\hline
	
	Tần suất sử dụng &
	Khi Student yêu cầu hoặc sau các mốc đánh giá quan trọng. \\
	\hline
	
\end{longtable}

\subsection{Luồng xử lý chính}

\begin{enumerate}
	\item Student mở chức năng dự đoán kết quả.
	
	\item Hệ thống kiểm tra Student đã thiết lập điểm mục tiêu.
	
	\item Hệ thống lấy hồ sơ năng lực mới nhất.
	
	\item Hệ thống lấy kết quả các bài đánh giá và thi thử.
	
	\item Hệ thống lấy tiến độ hoàn thành lộ trình.
	
	\item Hệ thống lấy tần suất và thời lượng luyện tập.
	
	\item Hệ thống xác định thời gian còn lại trước ngày thi.
	
	\item Hệ thống kiểm tra mức độ đầy đủ của dữ liệu.
	
	\item Hệ thống chuẩn hóa dữ liệu đầu vào.
	
	\item Thành phần dự đoán xử lý dữ liệu.
	
	\item Hệ thống tính xác suất đạt điểm mục tiêu.
	
	\item Hệ thống ước lượng khoảng điểm dự kiến.
	
	\item Hệ thống xác định các yếu tố ảnh hưởng chính.
	
	\item Hệ thống tạo gợi ý cải thiện.
	
	\item Hệ thống hiển thị kết quả và mức độ tin cậy.
	
	\item Hệ thống lưu phiên bản dự đoán để so sánh trong tương lai.
\end{enumerate}

\subsection{Luồng thay thế và trường hợp lỗi}

\begin{enumerate}[label=\textbf{A\arabic*.}]
	\item \textbf{Student chưa thiết lập điểm mục tiêu}
	
	Hệ thống yêu cầu Student thiết lập điểm mục tiêu trước khi thực hiện
	dự đoán.
	
	\item \textbf{Chưa đủ dữ liệu}
	
	Hệ thống thông báo chưa thể dự đoán đáng tin cậy và đề xuất Student
	thực hiện thêm bài đánh giá hoặc thi thử.
	
	\item \textbf{Dữ liệu có giá trị bất thường}
	
	Hệ thống loại bỏ hoặc đánh dấu dữ liệu bất thường trước khi dự đoán.
	
	\item \textbf{Thành phần dự đoán không hoạt động}
	
	Hệ thống thông báo chức năng tạm thời chưa khả dụng.
	
	\item \textbf{Mức độ tin cậy thấp}
	
	Hệ thống vẫn hiển thị kết quả nhưng kèm cảnh báo rõ ràng.
	
	\item \textbf{Điểm mục tiêu không phù hợp với dữ liệu hiện tại}
	
	Hệ thống không thay đổi mục tiêu của Student nhưng đưa ra cảnh báo và
	gợi ý kế hoạch cải thiện.
\end{enumerate}

\subsection{Dữ liệu đầu vào và kết quả đầu ra}

\begin{longtable}{|p{3.5cm}|p{5.3cm}|p{5.6cm}|}
	\caption{Dữ liệu của UC-17}
	\label{tab:uc17-data}\\
	\hline
	\textbf{Nhóm dữ liệu} &
	\textbf{Thành phần} &
	\textbf{Mô tả} \\
	\hline
	\endfirsthead
	
	\hline
	\textbf{Nhóm dữ liệu} &
	\textbf{Thành phần} &
	\textbf{Mô tả} \\
	\hline
	\endhead
	
	Đầu vào &
	Hồ sơ năng lực &
	Mức độ thành thạo hiện tại theo từng năng lực. \\
	\hline
	
	Đầu vào &
	Kết quả thi thử &
	Điểm số, thời gian làm và xu hướng qua các lần thi. \\
	\hline
	
	Đầu vào &
	Tiến độ lộ trình &
	Tỷ lệ nhiệm vụ đã hoàn thành và mức độ đúng hạn. \\
	\hline
	
	Đầu vào &
	Hoạt động luyện tập &
	Tần suất, thời lượng và kết quả luyện tập. \\
	\hline
	
	Đầu vào &
	Điểm mục tiêu &
	Điểm Student muốn đạt được. \\
	\hline
	
	Đầu vào &
	Thời gian còn lại &
	Số ngày hoặc tuần trước thời điểm thi. \\
	\hline
	
	Đầu ra &
	Xác suất đạt mục tiêu &
	Tỷ lệ ước lượng khả năng Student đạt điểm đã đặt. \\
	\hline
	
	Đầu ra &
	Khoảng điểm dự kiến &
	Khoảng điểm mà Student có khả năng đạt được. \\
	\hline
	
	Đầu ra &
	Mức độ tin cậy &
	Thông tin thể hiện mức độ chắc chắn của dự đoán. \\
	\hline
	
	Đầu ra &
	Yếu tố ảnh hưởng &
	Các năng lực, kết quả hoặc hành vi tác động nhiều đến dự đoán. \\
	\hline
	
	Đầu ra &
	Gợi ý cải thiện &
	Nội dung Student cần ưu tiên để tăng khả năng đạt mục tiêu. \\
	\hline
	
\end{longtable}

\subsection{Mối liên hệ với hồ sơ năng lực và lộ trình}

UC-17 sử dụng trực tiếp dữ liệu từ hồ sơ năng lực và tiến độ lộ trình.
Khi kết quả dự đoán thấp hơn yêu cầu, hệ thống gửi thông tin cho UC-11
để đề xuất tăng mức độ ưu tiên của các chủ đề còn yếu.

Kết quả dự đoán không trực tiếp sửa hồ sơ năng lực. Hồ sơ năng lực chỉ
được cập nhật từ dữ liệu học tập thực tế như đánh giá, thi thử và luyện
tập.

% =========================================================
\section{Mối quan hệ giữa các chức năng AI}
\label{sec:ai-function-relationship}
% =========================================================

Sáu chức năng AI và cá nhân hóa có quan hệ dữ liệu liên tục với nhau:

\begin{enumerate}
	\item UC-08 tạo kết quả chấm điểm và phân tích bài làm.
	
	\item UC-09 sử dụng kết quả đó để cập nhật hồ sơ năng lực.
	
	\item UC-11 sử dụng hồ sơ năng lực để tạo lộ trình học cá nhân hóa.
	
	\item UC-12 lựa chọn bài luyện tập dựa trên năng lực và lộ trình.
	
	\item Kết quả UC-12 tiếp tục quay lại cập nhật UC-09.
	
	\item UC-14 hỗ trợ Student trong quá trình thực hiện lộ trình và luyện
	tập.
	
	\item UC-17 sử dụng hồ sơ năng lực, tiến độ lộ trình và kết quả luyện
	tập để dự đoán khả năng đạt điểm mục tiêu.
	
	\item Kết quả dự đoán có thể tạo đề xuất điều chỉnh cho UC-11.
\end{enumerate}

Luồng dữ liệu tổng quát được mô tả như sau:

\begin{center}
	UC-08 Chấm điểm
	$\rightarrow$
	UC-09 Hồ sơ năng lực
	$\rightarrow$
	UC-11 Lộ trình học
	$\rightarrow$
	UC-12 Luyện tập thích ứng
	$\rightarrow$
	UC-09 Hồ sơ năng lực
\end{center}

\begin{center}
	UC-09 Hồ sơ năng lực
	$+$
	UC-11 Tiến độ lộ trình
	$\rightarrow$
	UC-17 Dự đoán điểm mục tiêu
\end{center}

AI Tutor hỗ trợ xuyên suốt quá trình Student học tập, luyện tập và xem
lộ trình. Toàn bộ các chức năng phải sử dụng thống nhất thông tin Student,
môn học, chủ đề, năng lực và lịch sử học tập.